\documentclass[12pt, a4paper]{book}

% ─── Core packages ────────────────────────────────────────────────────────────
\usepackage[T1]{fontenc}
\usepackage[utf8]{inputenc}
\usepackage[charter, cal=cmcal]{mathdesign}
\usepackage[scaled=0.92]{helvet}
\usepackage{geometry}
\geometry{
  a4paper,
  top=30mm, bottom=30mm,
  inner=32mm, outer=25mm,
  marginparwidth=18mm
}

% ─── Mathematics ──────────────────────────────────────────────────────────────
\usepackage{amsmath, amsthm}

% ─── Chemistry ────────────────────────────────────────────────────────────────
\usepackage[version=4]{mhchem}

% ─── Graphics & colour ────────────────────────────────────────────────────────
\usepackage{graphicx}
\usepackage[dvipsnames,svgnames,table]{xcolor}
\usepackage{tikz}
\usetikzlibrary{arrows.meta, shapes.geometric, positioning,
                calc, backgrounds, decorations.pathmorphing,
                decorations.markings, patterns, fadings}
\usepackage{pgfplots}
\pgfplotsset{compat=1.18}
\usepackage{tcolorbox}
\tcbuselibrary{skins, breakable, theorems}

% ─── Layout ───────────────────────────────────────────────────────────────────
\usepackage{titlesec}
\usepackage{fancyhdr}
\usepackage{enumitem}
\usepackage{booktabs}
\usepackage{array}
\usepackage{multicol}
\usepackage{marginnote}
\usepackage{caption}
\usepackage{parskip}
\usepackage{siunitx}
\usepackage{pifont}
\usepackage[colorlinks=true, linkcolor=BrickRed,
            urlcolor=NavyBlue, citecolor=ForestGreen]{hyperref}

% ══════════════════════════════════════════════════════════════════════════════
% COLOUR PALETTE
% ══════════════════════════════════════════════════════════════════════════════
\definecolor{JourneyColor}{HTML}{4A235A}   % deep purple for Journey 2
\definecolor{ModuleColor}{HTML}{6C3483}    % medium purple
\definecolor{PuzzleColor}{HTML}{1A5276}    % dark blue
\definecolor{DiscoveryColor}{HTML}{B7950B} % amber
\definecolor{WarningColor}{HTML}{C0392B}   % red
\definecolor{ProblemColor}{HTML}{117A65}   % teal
\definecolor{ReflectColor}{HTML}{1D6A39}   % green
\definecolor{HistoryColor}{HTML}{6E2F1A}   % brown
\definecolor{LightBg}{HTML}{F8F9FA}
\definecolor{ElectronBlue}{HTML}{1565C0}   % for electron diagrams
\definecolor{FieldRed}{HTML}{B71C1C}       % for field lines

% ══════════════════════════════════════════════════════════════════════════════
% TCOLORBOX ENVIRONMENTS
% ══════════════════════════════════════════════════════════════════════════════
\newtcolorbox{puzzlebox}[1][]{
  enhanced, breakable,
  colback=PuzzleColor!6, colframe=PuzzleColor,
  fonttitle=\bfseries\sffamily, title={The Puzzle},
  attach boxed title to top left={yshift=-2mm, xshift=6mm},
  boxed title style={colback=PuzzleColor, colframe=PuzzleColor},
  borderline west={3pt}{0pt}{PuzzleColor}, #1}

\newtcolorbox{discoverybox}[1][]{
  enhanced, breakable,
  colback=DiscoveryColor!8, colframe=DiscoveryColor,
  fonttitle=\bfseries\sffamily, title={Discovery},
  attach boxed title to top left={yshift=-2mm, xshift=6mm},
  boxed title style={colback=DiscoveryColor, colframe=DiscoveryColor},
  borderline west={3pt}{0pt}{DiscoveryColor}, #1}

\newtcolorbox{problembox}[1][]{
  enhanced, breakable,
  colback=ProblemColor!5, colframe=ProblemColor,
  fonttitle=\bfseries\sffamily, title={Problems},
  attach boxed title to top left={yshift=-2mm, xshift=6mm},
  boxed title style={colback=ProblemColor, colframe=ProblemColor},
  borderline west={3pt}{0pt}{ProblemColor}, #1}

\newtcolorbox{deepbox}[1][]{
  enhanced, breakable,
  colback=JourneyColor!5, colframe=JourneyColor,
  fonttitle=\bfseries\sffamily, title={Deep Problems},
  attach boxed title to top left={yshift=-2mm, xshift=6mm},
  boxed title style={colback=JourneyColor, colframe=JourneyColor},
  borderline west={4pt}{0pt}{JourneyColor}, #1}

\newtcolorbox{historybox}[1][]{
  enhanced, breakable,
  colback=HistoryColor!7, colframe=HistoryColor!60,
  fonttitle=\bfseries\sffamily\small, title={Historical Note},
  attach boxed title to top left={yshift=-2mm, xshift=6mm},
  boxed title style={colback=HistoryColor!80, colframe=HistoryColor}, #1}

\newtcolorbox{reflectionbox}[1][]{
  enhanced, breakable,
  colback=ReflectColor!6, colframe=ReflectColor,
  fonttitle=\bfseries\sffamily, title={Reflection},
  attach boxed title to top left={yshift=-2mm, xshift=6mm},
  boxed title style={colback=ReflectColor, colframe=ReflectColor},
  borderline west={3pt}{0pt}{ReflectColor}, #1}

\newtcolorbox{modelbreaks}[1][]{
  enhanced, breakable,
  colback=WarningColor!6, colframe=WarningColor,
  fonttitle=\bfseries\sffamily, title={Where This Model Breaks},
  attach boxed title to top left={yshift=-2mm, xshift=6mm},
  boxed title style={colback=WarningColor, colframe=WarningColor},
  borderline west={4pt}{0pt}{WarningColor}, #1}

\newtcolorbox{namedconcept}[2][]{
  enhanced,
  colback=ModuleColor!6, colframe=ModuleColor,
  fonttitle=\bfseries\sffamily\large, title={#2},
  attach boxed title to top left={yshift=-2mm, xshift=6mm},
  boxed title style={colback=ModuleColor, colframe=ModuleColor}, #1}

% ══════════════════════════════════════════════════════════════════════════════
% THEOREM ENVIRONMENTS
% ══════════════════════════════════════════════════════════════════════════════
\theoremstyle{definition}
\newtheorem{observation}{Observation}[chapter]
\newtheorem{experiment}{Experiment}[chapter]

% ══════════════════════════════════════════════════════════════════════════════
% SECTION FORMATTING
% ══════════════════════════════════════════════════════════════════════════════
\titleformat{\chapter}[display]
  {\normalfont\sffamily\huge\bfseries\color{JourneyColor}}
  {\textcolor{JourneyColor!60}{\chaptertitlename\ \thechapter}}
  {10pt}{\Huge}[\vspace{2pt}\textcolor{JourneyColor}{\titlerule[1.5pt]}]

\titleformat{\section}
  {\normalfont\sffamily\Large\bfseries\color{ModuleColor}}
  {\thesection}{1em}{}[\textcolor{ModuleColor!40}{\titlerule[0.6pt]}]

\titleformat{\subsection}
  {\normalfont\sffamily\large\bfseries\color{JourneyColor!80}}
  {\thesubsection}{1em}{}

\titleformat{\subsubsection}
  {\normalfont\sffamily\normalsize\bfseries}
  {\thesubsubsection}{1em}{}

% ══════════════════════════════════════════════════════════════════════════════
% HEADER / FOOTER
% ══════════════════════════════════════════════════════════════════════════════
\pagestyle{fancy}
\fancyhf{}
\fancyhead[LE]{\small\sffamily\textcolor{JourneyColor}{\leftmark}}
\fancyhead[RO]{\small\sffamily\textcolor{ModuleColor}{\rightmark}}
\fancyfoot[C]{\small\thepage}
\renewcommand{\headrulewidth}{0.4pt}

% ══════════════════════════════════════════════════════════════════════════════
% CUSTOM COMMANDS
% ══════════════════════════════════════════════════════════════════════════════
\newcommand{\chem}[1]{\ce{#1}}
\newcommand{\important}[1]{\textcolor{ModuleColor}{\textbf{#1}}}
\newcommand{\newterm}[1]{\textit{\textcolor{DiscoveryColor}{#1}}}
\newcommand{\electron}{$e^-$}
\newcommand{\proton}{$p^+$}

% ══════════════════════════════════════════════════════════════════════════════
% DOCUMENT
% ══════════════════════════════════════════════════════════════════════════════
\begin{document}
\setcounter{chapter}{1}   % Journey 2 = Chapter 2

% ─── Journey 2 opening ────────────────────────────────────────────────────────
\chapter[Journey 2 --- Architecture of the Atom]
        {Journey Two\\[8pt]
         \textcolor{ModuleColor}{\large Architecture of the Atom}}
\label{ch:journey2}

\begin{tcolorbox}[
  enhanced, colback=JourneyColor!8, colframe=JourneyColor,
  arc=4pt, left=6mm, right=6mm, top=4mm, bottom=4mm
]
\sffamily\small
\textbf{What this journey is about.}

Journey~1 ended with a triumph and a mystery.
The triumph: matter is made of atoms, each element has one type,
and the evidence converges from five independent directions.
The mystery: we treated atoms as featureless spheres.
We knew their relative masses, their ratios in compounds,
and that they could bond together.
We knew nothing about what was \emph{inside} them.

This journey descends into the atom itself.
What we find there is far stranger than anyone expected.
The atom is almost entirely empty space.
It contains a tiny, dense, positively charged nucleus
surrounded by a cloud of negatively charged electrons.
Those electrons are not planets in fixed orbits.
They are waves of probability.
Their energy is not continuous but comes in discrete packets.

None of this was assumed.
All of it was \emph{forced} on physicists and chemists
by experimental evidence they could not explain any other way.

We will follow the evidence exactly as it arrived:
first the electrons, then the nucleus, then the catastrophe
of classical physics, then the quantum revolution.
By the end, you will understand not just what atoms look like
but \emph{why} they must look that way.
\end{tcolorbox}

\vspace{1em}

% Journey 2 concept flow
\begin{center}
\begin{tikzpicture}[
  every node/.style={font=\sffamily\footnotesize, align=center},
  box/.style={draw=JourneyColor!50, fill=JourneyColor!7,
              rounded corners=3pt, minimum width=2.8cm,
              minimum height=0.8cm},
  arr/.style={-{Stealth[length=4pt]}, thick, color=ModuleColor!70}
]
  \node[box] (e)   at (0,   0) {Electrons\\discovered (2.1)};
  \node[box] (nuc) at (3.2, 0) {Nucleus\\discovered (2.2)};
  \node[box] (cat) at (6.4, 0) {Classical physics\\fails (2.3)};
  \node[box] (spec)at (9.6, 0) {Spectra demand\\quanta (2.4)};
  \node[box] (bohr)at (0,  -1.8){Bohr model\\(2.5)};
  \node[box] (wave)at (3.2,-1.8){Wave-particle\\duality (2.6)};
  \node[box] (orb) at (6.4,-1.8){Orbitals and\\quantum numbers (2.7)};
  \node[box, draw=DiscoveryColor, fill=DiscoveryColor!10,
        minimum width=3.5cm]
             (per) at (9.6,-1.8){Periodic table\\explained (J3)};

  \draw[arr] (e) -- (nuc);
  \draw[arr] (nuc) -- (cat);
  \draw[arr] (cat) -- (spec);
  \draw[arr] (spec) -- (bohr);
  \draw[arr] (bohr) -- (wave);
  \draw[arr] (wave) -- (orb);
  \draw[arr] (orb)  -- (per);
  \draw[arr, dashed] (spec.south) -- ++(0,-0.4) -- ++(-9.6,0) -- (bohr.north);
\end{tikzpicture}
\end{center}

\vspace{1.5em}

% ══════════════════════════════════════════════════════════════════════════════
%  MODULE 2.1 — THE ATOM GETS STRUCTURE: THOMSON AND THE ELECTRON
% ══════════════════════════════════════════════════════════════════════════════
\section{Module 2.1 --- The Atom Gets Structure: Thomson and the Electron}
\label{sec:mod2.1}

\subsection*{Motivation}

Dalton's atom had no internal structure.
It was a tiny, solid, indivisible sphere.
The only property that distinguished one element's atoms
from another's was mass.

But chemistry kept producing puzzles that Dalton's model
could not address.
Why does electricity decompose compounds?
Why does rubbing amber attract feathers?
Why does passing electricity through gases produce light?

These questions all involve \emph{charged particles}.
Something inside the atom must carry charge.
To find it, physicists turned to a peculiar piece of apparatus:
the cathode ray tube.

% ─── Section: The Cathode Ray Tube ───────────────────────────────────────────
\subsection{The Cathode Ray Tube: A Window into the Atom}

A cathode ray tube is a glass tube from which most of the air
has been pumped out.
Metal electrodes are sealed into each end.
When a high voltage is applied, something remarkable happens:
a faint greenish glow appears, and if a screen is placed
inside the tube, it casts a sharp shadow.

Something is travelling in straight lines from the
negative electrode (cathode) to the positive electrode (anode).
This something was called a \newterm{cathode ray}.

\begin{center}
\begin{tikzpicture}[
  every node/.style={font=\sffamily\small},
  scale=1.0
]
  % Tube outline
  \draw[thick, JourneyColor!50, rounded corners=6pt]
    (-5.5,-1.2) rectangle (5.5, 1.2);

  % Cathode (left electrode)
  \fill[JourneyColor!40] (-5.0,-0.8) rectangle (-4.6, 0.8);
  \node[font=\sffamily\footnotesize, text=JourneyColor]
    at (-4.8,-1.05) {Cathode $-$};

  % Anode (right electrode, with hole)
  \fill[JourneyColor!40] (4.4,-0.8) rectangle (4.8, 0.8);
  \fill[white] (4.4,-0.15) rectangle (4.8, 0.15);
  \node[font=\sffamily\footnotesize, text=JourneyColor]
    at (4.6,-1.05) {Anode $+$};

  % Cathode ray beam
  \draw[-{Stealth[length=6pt]}, ultra thick, ElectronBlue!70,
        line width=2pt]
    (-4.5, 0) -- (4.3, 0);
  \node[font=\sffamily\footnotesize, text=ElectronBlue]
    at (-0.5, 0.35) {cathode ray};

  % Glowing screen at right end
  \fill[DiscoveryColor!30] (4.9,-1.1) rectangle (5.4, 1.1);
  \fill[DiscoveryColor!80, opacity=0.6] (5.1,-0.3) circle (0.25cm);
  \node[font=\sffamily\footnotesize, text=HistoryColor]
    at (5.2,-1.0) {glow};

  % Voltage source
  \draw[thick, JourneyColor!50]
    (-4.8,-1.2) -- (-4.8,-1.7) -- (4.6,-1.7) -- (4.6,-1.2);
  \node[draw=JourneyColor!50, fill=LightBg, rounded corners=2pt,
        font=\sffamily\footnotesize]
    at (0,-1.7) {high voltage};

  % Vacuum annotation
  \node[font=\sffamily\footnotesize\itshape, text=JourneyColor!60]
    at (0, 0.85) {evacuated glass tube};
\end{tikzpicture}
\end{center}

\vspace{0.5em}

The cathode ray had several strange properties that experimenters
catalogued over decades:

\begin{itemize}[leftmargin=2em]
  \item It always travels from cathode (negative) to anode (positive),
        regardless of which metal the electrodes are made of.
  \item It is deflected by magnetic fields, just as a current-carrying
        wire is.
  \item It is deflected toward a positive plate and away from a
        negative plate --- suggesting it carries \emph{negative charge}.
  \item It can push a small paddle wheel placed in its path,
        suggesting it has momentum and therefore mass.
  \item Its properties are \emph{identical} regardless of which
        metal the cathode is made of, or which gas is in the tube.
\end{itemize}

The last point is crucial.
If the ray came from the metal itself, different metals
should produce different rays.
The fact that all metals produce the same ray suggests
the ray comes from something \emph{universal} ---
something present in all matter.

\begin{puzzlebox}
Before Thomson's interpretation, several hypotheses existed
for what cathode rays actually were:

\begin{enumerate}[label=(\alph*), leftmargin=*]
  \item \textbf{Hypothesis A: Cathode rays are electromagnetic waves}
        (like light), just at shorter wavelength.
        \textit{Evidence for:} they travel in straight lines,
        pass through thin foils.
        \textit{Evidence against:} they are deflected by magnetic
        fields (light is not); they push paddle wheels (light
        pressure is negligible).

  \item \textbf{Hypothesis B: Cathode rays are streams of charged particles.}
        \textit{Evidence for:} deflected by electric and magnetic fields;
        carry momentum; same charge sign regardless of cathode material.
        \textit{Evidence against:} none, once you accept the experiment.

  Consider the deflection by electric fields.
  A positive plate is placed above the beam, a negative plate below.

  \begin{enumerate}[label=(\roman*)]
    \item Which way would the beam deflect if cathode rays
          carry negative charge?
    \item Which way if they carry positive charge?
    \item Thomson observed deflection \emph{toward the positive plate}.
          Which hypothesis does this support?
    \item If cathode rays are electromagnetic waves (hypothesis A),
          what would you expect when you place an electric field
          across them?
  \end{enumerate}
\end{enumerate}
\end{puzzlebox}

% ─── Section: Thomson's Measurement ──────────────────────────────────────────
\subsection{Thomson's Crucial Measurement: The Charge-to-Mass Ratio}

In 1897, J.~J.~Thomson went beyond simply demonstrating that
cathode rays are charged particles.
He measured a specific property of those particles:
the ratio of their charge to their mass, $e/m$.

His method was elegant.
He applied simultaneous electric and magnetic fields to the beam
and adjusted them until the deflections cancelled ---
meaning the electric and magnetic forces were equal and opposite.

\begin{center}
\begin{tikzpicture}[
  every node/.style={font=\sffamily\small},
  scale=1.0
]
  % Tube outline
  \draw[thick, JourneyColor!40, rounded corners=4pt]
    (-5,-1.4) rectangle (5.5, 1.4);

  % Cathode
  \fill[JourneyColor!40] (-4.8,-0.9) rectangle (-4.4, 0.9);
  \node[font=\sffamily\footnotesize] at (-4.6,-1.2) {$-$};

  % Anode with slit
  \fill[JourneyColor!40] (-3.2,-0.9) rectangle (-2.8, 0.9);
  \fill[white] (-3.2,-0.15) rectangle (-2.8, 0.15);

  % Deflection plates (electric field)
  \fill[FieldRed!30] (-1.8, 0.6) rectangle (1.8, 0.9);
  \fill[ElectronBlue!30] (-1.8,-0.9) rectangle (1.8,-0.6);
  \node[font=\sffamily\footnotesize, text=FieldRed]
    at (0, 1.05) {$+$ plate (electric)};
  \node[font=\sffamily\footnotesize, text=ElectronBlue]
    at (0,-1.1) {$-$ plate (electric)};

  % Magnetic field arrows (into page)
  \foreach \x in {-1.4,-0.6,0.2,1.0} {
    \node[font=\tiny, text=ModuleColor!70] at (\x, 0) {$\otimes$};
  }
  \node[font=\sffamily\footnotesize, text=ModuleColor]
    at (3.2, 0.35) {$\vec{B}$ into page};

  % Beam - straight (balanced forces)
  \draw[-{Stealth[length=6pt]}, ultra thick, ElectronBlue!80,
        line width=2pt]
    (-4.3,0) -- (4.8, 0);
  \node[font=\sffamily\footnotesize, text=ElectronBlue]
    at (3.5,-0.35) {beam undeflected};

  % Force annotations
  \draw[-{Stealth}, thick, FieldRed]
    (0.5, 0.05) -- (0.5, 0.55)
    node[right, font=\sffamily\footnotesize, text=FieldRed]{$F_E = eE$};
  \draw[-{Stealth}, thick, ElectronBlue!80]
    (-0.5, 0.05) -- (-0.5,-0.45)
    node[right, font=\sffamily\footnotesize, text=ElectronBlue]{$F_B = evB$};

  % Screen
  \fill[DiscoveryColor!25] (4.9,-1.3) rectangle (5.4, 1.3);
  \fill[DiscoveryColor!70] (5.0,-0.15) circle (0.2cm);
\end{tikzpicture}
\end{center}

\vspace{0.5em}

When the forces balance:
\[
  eE = evB \quad \Rightarrow \quad v = \frac{E}{B}
\]

Thomson then turned off the electric field and let the magnetic
field alone deflect the beam in a circular arc of radius $r$.
For circular motion:
\[
  evB = \frac{mv^2}{r} \quad \Rightarrow \quad
  \frac{e}{m} = \frac{v}{rB} = \frac{E}{rB^2}
\]

Every measurable quantity on the right side ($E$, $r$, $B$)
was under Thomson's control.
By measuring the radius of curvature and knowing the fields,
he could calculate $e/m$ directly.

His result:
\[
  \frac{e}{m} \approx 1.76 \times 10^{11}\ \si{\coulomb\per\kilogram}
\]

This number is extraordinary.
The largest $e/m$ ratio previously known was for
the hydrogen ion in electrolysis:
approximately $9.6 \times 10^7\ \si{\coulomb\per\kilogram}$.

The cathode ray particle had an $e/m$ ratio
\emph{about 1800 times larger} than the hydrogen ion.

\begin{puzzlebox}
Thomson measured $e/m \approx 1.76 \times 10^{11}\ \si{\coulomb\per\kilogram}$
for cathode ray particles,
compared to $e/m \approx 9.6 \times 10^7\ \si{\coulomb\per\kilogram}$
for the hydrogen ion.

\begin{enumerate}[label=(\alph*), leftmargin=*]
  \item Calculate the ratio of these two $e/m$ values.
        By what factor does the cathode ray particle's ratio
        exceed that of the hydrogen ion?

  \item Two interpretations are possible:
        \begin{itemize}[nosep]
          \item The cathode ray particle has the \emph{same charge}
                as the hydrogen ion but much smaller mass.
          \item The cathode ray particle has the \emph{same mass}
                as the hydrogen ion but much larger charge.
        \end{itemize}
        Thomson argued the first interpretation was correct.
        What additional experimental evidence would you need
        to distinguish between them?

  \item Millikan later measured the charge of the electron
        directly as $e = 1.602 \times 10^{-19}\ \si{\coulomb}$.
        Use this and Thomson's $e/m$ ratio to calculate
        the mass of the electron.

  \item The mass of a hydrogen atom is
        $1.67 \times 10^{-27}\ \si{\kilogram}$.
        What fraction of the hydrogen atom's mass is the electron?
        What does this tell you about where the atom's mass is?
\end{enumerate}
\end{puzzlebox}

\vspace{0.5em}

Thomson's measurement established something profound:
there exists a particle, present in all matter,
with a mass roughly 1/1836 that of the lightest atom.
He called it the \newterm{electron}.

\begin{namedconcept}{The Electron (Thomson, 1897)}
The \newterm{electron} is a fundamental particle with:
\begin{itemize}[nosep, leftmargin=2em]
  \item Mass: $m_e = 9.109 \times 10^{-31}\ \si{\kilogram}$
  \item Charge: $e = -1.602 \times 10^{-19}\ \si{\coulomb}$ (negative)
  \item Present in all matter, independent of the element
  \item About 1/1836 the mass of a hydrogen atom
\end{itemize}
Electrons are the carriers of electric current.
They are the particles that are shared or transferred
when atoms bond together.
Their behaviour in atoms determines the
chemical properties of every element.
\end{namedconcept}

% ─── Section: The Plum Pudding Model ─────────────────────────────────────────
\subsection{The First Atomic Model: Plum Pudding}

Atoms are electrically neutral.
Electrons are negatively charged.
Therefore something in the atom must carry positive charge,
and its positive charge must exactly cancel the negative
charge of the electrons.

Thomson proposed a model:
the atom is a sphere of uniformly distributed positive charge,
like a pudding, with electrons embedded in it like plums
(or raisins --- the British and American analogies differ).

\begin{center}
\begin{tikzpicture}[
  every node/.style={font=\sffamily\small},
  scale=1.1
]
  % Positive sphere
  \shade[ball color=FieldRed!15] (0,0) circle (1.8cm);
  \draw[thick, FieldRed!40] (0,0) circle (1.8cm);
  \node[font=\sffamily\footnotesize, text=FieldRed!70]
    at (0,-0.3) {uniform positive charge};

  % Electrons as dots
  \foreach \x/\y in {
    0.7/1.0, -0.9/0.8, 1.1/-0.4,
    -1.2/-0.5, 0.2/-1.2, -0.4/0.4,
    1.3/0.3, -0.6/-1.1
  }{
    \fill[ElectronBlue] (\x,\y) circle (0.12cm);
    \draw[ElectronBlue!60, thick] (\x,\y) circle (0.12cm);
  }

  % Legend
  \fill[ElectronBlue] (3.2, 0.5) circle (0.12cm);
  \node[right, font=\sffamily\footnotesize, text=ElectronBlue]
    at (3.35, 0.5) {electron ($e^-$)};
  \fill[FieldRed!20] (3.2, -0.1) circle (0.3cm);
  \draw[FieldRed!50, thick] (3.2,-0.1) circle (0.3cm);
  \node[right, font=\sffamily\footnotesize, text=FieldRed]
    at (3.55,-0.1) {positive charge cloud};

  % Title
  \node[font=\sffamily\small\bfseries, color=JourneyColor]
    at (0, 2.3) {Thomson's ``Plum Pudding'' Model};
\end{tikzpicture}
\end{center}

\vspace{0.5em}

The plum pudding model made specific predictions.
Most importantly: since the positive charge is spread uniformly
through the atom, any particle fired at the atom should
pass through it, deflected only slightly by the weak,
diffuse field.

There should be \emph{no large-angle deflections}.

This prediction was about to be tested.
And it would be spectacularly wrong.

\begin{historybox}
\textbf{Thomson's Other Legacy: The Discovery of Isotopes}

Thomson's cathode ray apparatus was later refined by Francis Aston
into the \emph{mass spectrometer}.
By deflecting ionised atoms in magnetic fields ---
using the same principle Thomson used for electrons ---
Aston could sort atoms by their mass-to-charge ratio.

In 1913, when Thomson passed ionised neon through his apparatus,
he found two distinct beams, corresponding to masses 20 and 22.
This was the first direct evidence for \emph{isotopes}:
atoms of the same element with different masses.

The weighted average (90\% mass-20, 10\% mass-22) gives
the observed atomic mass of $20.2$ --- exactly what was
measured by chemical methods.

This confirmed Dalton's anomaly: atoms of the same element
are \emph{not} all identical in mass.
\end{historybox}

% ─── Section: Electrolysis and Charge Quantisation ───────────────────────────
\subsection{Faraday's Earlier Clue: Charge Comes in Packets}

Thomson's discovery of the electron did not come entirely
from nowhere.
Fifty years earlier, Michael Faraday had studied electrolysis
and noticed something that, in hindsight, screamed
``discrete charge carriers.''

Faraday found two laws of electrolysis (1833):

\begin{namedconcept}{Faraday's Laws of Electrolysis}
\textbf{First law:} The mass of substance deposited at an electrode
is proportional to the total charge passed.

\textbf{Second law:} The masses of different substances deposited
by the same charge are proportional to their
molar masses divided by a small integer
(their \emph{valence}).

\[
  m = \frac{Q \cdot M}{n_{\text{val}} \cdot F}
\]

where $Q$ is charge, $M$ is molar mass,
$n_{\text{val}}$ is valence (1, 2, or 3),
and $F = \SI{96485}{\coulomb\per\mol}$ is the Faraday constant.
\end{namedconcept}

\vspace{0.5em}

The Faraday constant represents the charge required to deposit
one mole of a singly-charged ion.
That it is a \emph{fixed, universal constant} is a strong
hint that charge comes in discrete packets ---
one packet per ion.

\begin{problembox}
\begin{enumerate}[label=\textbf{\arabic*.}]
  \item A current of \SI{2.00}{\ampere} passes through a solution
        of copper sulfate (\ce{CuSO4}) for \SI{30.0}{\minute}.
        Copper has molar mass \SI{63.5}{\gram\per\mol} and
        valence 2.
        \begin{enumerate}[label=(\alph*)]
          \item Calculate the total charge passed in coulombs.
          \item Calculate the mass of copper deposited.
          \item How many moles of copper were deposited?
          \item How many copper atoms were deposited?
        \end{enumerate}

  \item Electrolysis of molten \ce{Al2O3} produces aluminium.
        Aluminium has molar mass \SI{27.0}{\gram\per\mol}
        and valence 3.
        What current must flow for \SI{2.00}{\hour}
        to deposit \SI{27.0}{\gram} of aluminium?

  \item The Faraday constant $F = \SI{96485}{\coulomb\per\mol}$
        can be interpreted as: one mole of electrons
        carries a charge of \SI{96485}{\coulomb}.
        \begin{enumerate}[label=(\alph*)]
          \item Using $e = 1.602\times10^{-19}\ \si{\coulomb}$,
                calculate the charge of one mole of electrons.
          \item How close is this to the Faraday constant?
          \item What does this tell us about the relationship
                between the Faraday constant and Avogadro's number?
        \end{enumerate}
\end{enumerate}
\end{problembox}

% ─── Section: Millikan's Oil Drop ────────────────────────────────────────────
\subsection{Millikan's Oil Drop: Measuring a Single Electron}

Thomson knew the \emph{ratio} $e/m$ but not $e$ or $m$ separately.
To find the actual charge on an electron,
Robert Millikan devised a beautifully simple experiment in 1909.

He sprayed fine oil droplets between two charged plates.
As the drops fell through the air, some acquired a charge
by friction --- picking up one, two, or a few electrons
from the surrounding air.

By adjusting the electric field between the plates,
Millikan could make a charged droplet hover:
the upward electric force exactly balancing gravity.

\begin{center}
\begin{tikzpicture}[
  every node/.style={font=\sffamily\small},
  scale=1.0
]
  % Plates
  \fill[FieldRed!20] (-3.5, 1.8) rectangle (3.5, 2.2);
  \fill[ElectronBlue!20] (-3.5,-1.8) rectangle (3.5,-2.2);
  \node[font=\sffamily\footnotesize, text=FieldRed]
    at (-4.2, 2.0) {$+$ plate};
  \node[font=\sffamily\footnotesize, text=ElectronBlue]
    at (-4.2,-2.0) {$-$ plate};

  % Electric field lines
  \foreach \x in {-2.5,-1.5,-0.5,0.5,1.5,2.5}{
    \draw[-{Stealth[length=3pt]}, thin, FieldRed!40]
      (\x, 1.7) -- (\x,-1.7);
  }
  \node[font=\sffamily\footnotesize, text=FieldRed!60]
    at (4.0, 0) {$\vec{E}$};

  % Oil drop
  \shade[ball color=HistoryColor!40] (0,0) circle (0.22cm);
  \node[font=\sffamily\footnotesize, text=HistoryColor]
    at (1.2, 0) {oil drop (charged)};

  % Force arrows
  \draw[-{Stealth[length=5pt]}, ultra thick, FieldRed, line width=1.5pt]
    (0, 0.25) -- (0, 1.1)
    node[right, font=\sffamily\footnotesize, text=FieldRed]{$F_E = qE$};
  \draw[-{Stealth[length=5pt]}, ultra thick, JourneyColor, line width=1.5pt]
    (0,-0.25) -- (0,-1.1)
    node[right, font=\sffamily\footnotesize, text=JourneyColor]{$F_g = mg$};

  % Equilibrium annotation
  \node[draw=ModuleColor!60, fill=ModuleColor!8, rounded corners=2pt,
        font=\sffamily\footnotesize, text=ModuleColor]
    at (0,-2.8) {hover condition: $qE = mg$};
\end{tikzpicture}
\end{center}

\vspace{0.5em}

At equilibrium: $qE = mg$, so $q = mg/E$.

Millikan could measure $m$ (from the terminal velocity
when the field was off) and knew $E$ (from the voltage
and plate separation).
So he could calculate $q$ --- the charge on the droplet.

When he tabulated hundreds of measurements,
a pattern emerged that was impossible to miss:
all values of $q$ were \emph{integer multiples} of
a single smallest value.

No droplet ever had charge $1.5e$ or $2.7e$.
Every droplet had charge $1e$, $2e$, $3e$, $4e$, \ldots

\begin{namedconcept}{Quantisation of Charge}
Electric charge is \newterm{quantised}:
it always comes in integer multiples of the fundamental unit

\[
  e = 1.602 \times 10^{-19}\ \si{\coulomb}
\]

This is the charge of a single electron (negative)
or proton (positive).
No isolated particle has ever been observed
with charge smaller than $e$.
\end{namedconcept}

\begin{problembox}
\begin{enumerate}[label=\textbf{\arabic*.}]
  \item Millikan measures the following charges on oil droplets:
        \begin{center}
        \begin{tabular}{cc}
        \toprule
        Droplet & Charge (\si{\coulomb}) \\
        \midrule
        A & $3.204 \times 10^{-19}$ \\
        B & $6.408 \times 10^{-19}$ \\
        C & $1.602 \times 10^{-19}$ \\
        D & $4.806 \times 10^{-19}$ \\
        E & $8.010 \times 10^{-19}$ \\
        \bottomrule
        \end{tabular}
        \end{center}
        \begin{enumerate}[label=(\alph*)]
          \item Find the largest value that divides all
                five charges to give integer results.
                This is the elementary charge $e$.
          \item How many electrons does each droplet carry?
          \item If a droplet has charge $2.4\times10^{-19}\ \si{\coulomb}$,
                is this consistent with charge quantisation?
                What would you conclude?
        \end{enumerate}

  \item A charged oil droplet of mass
        $2.0 \times 10^{-15}\ \si{\kilogram}$
        hovers between plates with electric field
        $E = \SI{4.9e4}{\volt\per\metre}$.
        \begin{enumerate}[label=(\alph*)]
          \item Calculate the charge on the droplet.
                (Use $g = \SI{9.8}{\metre\per\second\squared}$.)
          \item How many electrons does the droplet carry?
          \item If the electric field is doubled,
                describe the subsequent motion of the droplet.
        \end{enumerate}
\end{enumerate}
\end{problembox}

% ─── Section: Where the Plum Pudding Model Stands ────────────────────────────
\subsection{The Model at the End of Module 2.1}

At the end of 1909, the atomic model looked like this:

\begin{center}
\begin{tikzpicture}[
  every node/.style={font=\sffamily\small, align=center}
]
  % Known
  \node[draw=ModuleColor!60, fill=ModuleColor!8,
        rounded corners=4pt, minimum width=4.5cm,
        minimum height=0.8cm]
        (k1) at (0, 2.0) {Atoms contain electrons};
  \node[draw=ModuleColor!60, fill=ModuleColor!8,
        rounded corners=4pt, minimum width=4.5cm,
        minimum height=0.8cm]
        (k2) at (0, 1.0) {Electron charge and mass known};
  \node[draw=ModuleColor!60, fill=ModuleColor!8,
        rounded corners=4pt, minimum width=4.5cm,
        minimum height=0.8cm]
        (k3) at (0, 0.0) {Charge is quantised};
  \node[draw=ModuleColor!60, fill=ModuleColor!8,
        rounded corners=4pt, minimum width=4.5cm,
        minimum height=0.8cm]
        (k4) at (0,-1.0) {Atom is electrically neutral};
  \node[font=\sffamily\small\bfseries, color=ModuleColor]
        at (0, 2.8) {Known};

  % Unknown
  \node[draw=WarningColor!60, fill=WarningColor!8,
        rounded corners=4pt, minimum width=4.5cm,
        minimum height=0.8cm]
        (u1) at (7.5, 2.0) {Where is positive charge?};
  \node[draw=WarningColor!60, fill=WarningColor!8,
        rounded corners=4pt, minimum width=4.5cm,
        minimum height=0.8cm]
        (u2) at (7.5, 1.0) {How is it distributed?};
  \node[draw=WarningColor!60, fill=WarningColor!8,
        rounded corners=4pt, minimum width=4.5cm,
        minimum height=0.8cm]
        (u3) at (7.5, 0.0) {Where are electrons inside?};
  \node[draw=WarningColor!60, fill=WarningColor!8,
        rounded corners=4pt, minimum width=4.5cm,
        minimum height=0.8cm]
        (u4) at (7.5,-1.0) {What holds atom together?};
  \node[font=\sffamily\small\bfseries, color=WarningColor]
        at (7.5, 2.8) {Unknown};

\end{tikzpicture}
\end{center}

\vspace{0.5em}

Thomson's answer to ``where is the positive charge?''
was the plum pudding model: spread uniformly throughout
a sphere the size of the atom.

This model is about to be tested --- and destroyed ---
by one of the most elegant experiments in the history of science.

\begin{modelbreaks}
Thomson's plum pudding model predicts that a fast particle
fired at an atom will pass straight through with only
\emph{slight} deflection, because the positive charge
is spread out and therefore weak at any one point.

\textbf{The test:} Rutherford fires alpha particles
(\ce{^4He^{2+}} nuclei, fast and heavy) at a thin gold foil.

\textbf{The prediction:} Almost all pass straight through.
Deflections of more than a few degrees should be essentially
impossible.

\textbf{What actually happened:} Module~2.2.
\end{modelbreaks}

% ─── Section: Deep Problems ───────────────────────────────────────────────────
\subsection*{Deep Problems}

\begin{deepbox}
\begin{enumerate}[label=\textbf{D\arabic*.}]

  \item \textbf{The Thomson Experiment in Full.}

        Thomson applies an electric field $E = \SI{2.0e4}{\volt\per\metre}$
        and a magnetic field $B = \SI{1.5e-3}{\tesla}$
        simultaneously, and the cathode ray beam is undeflected.

        \begin{enumerate}[label=(\alph*)]
          \item Calculate the speed of the electrons.
          \item He then turns off the electric field and measures
                a circular radius of curvature $r = \SI{0.12}{\metre}$.
                Calculate $e/m$.
          \item Compare your answer to the modern value
                $e/m = 1.759 \times 10^{11}\ \si{\coulomb\per\kilogram}$.
                What is the percentage error?
          \item In practice, Thomson's early measurements
                gave $e/m$ values that scattered considerably.
                Suggest two experimental sources of uncertainty
                that could cause this scatter.
        \end{enumerate}

  \item \textbf{Why Electrons Must Be Universal.}

        Thomson found that the $e/m$ ratio of cathode ray particles
        was identical regardless of:
        \begin{itemize}[nosep]
          \item the cathode metal (iron, platinum, aluminium)
          \item the residual gas in the tube
          \item the magnitude of the applied voltage
        \end{itemize}

        \begin{enumerate}[label=(\alph*)]
          \item Why does the independence from cathode material
                suggest that electrons are a \emph{universal constituent}
                of matter, not a product of the specific metal?
          \item Why does the independence from applied voltage
                matter? What would a varying $e/m$ with voltage imply?
          \item Construct the strongest possible argument, using
                only pre-1897 knowledge, that atoms must contain
                something smaller than themselves that carries charge.
                (Hints: electrolysis, sparks in gases, the production
                of light when metals are heated.)
          \item Today we know quarks have charges of $\pm e/3$
                and $\pm 2e/3$.
                Does this contradict Millikan's finding that
                charge is quantised in units of $e$?
                Explain carefully.
        \end{enumerate}

  \item \textbf{Forces Inside the Plum Pudding Atom.}

        Thomson proposed that electrons vibrate about equilibrium
        positions inside the positive sphere.
        For a uniform positive sphere of radius $R$ and total
        charge $+Ze$, the electric field \emph{inside} the sphere
        at distance $r$ from the centre is:
        \[
          E(r) = \frac{Ze}{4\pi\varepsilon_0}\cdot\frac{r}{R^3}
        \]
        (This follows from Gauss's law, which you may look up.)

        \begin{enumerate}[label=(\alph*)]
          \item Show that a single electron placed at distance
                $r$ from the centre of the sphere experiences
                a restoring force $F = -kr$ where $k$ is a constant.
                Find $k$ in terms of the given quantities.
          \item This is the equation of simple harmonic motion.
                What is the angular frequency $\omega$ of oscillation?
                (Recall $\omega = \sqrt{k/m_e}$.)
          \item For hydrogen ($Z=1$, $R \approx 5.3\times10^{-11}$ m),
                calculate the frequency of oscillation in Hz.
          \item What is the wavelength of light with this frequency?
                In what part of the electromagnetic spectrum is it?
          \item Does this match the observed hydrogen spectrum?
                (Hydrogen emits light at specific wavelengths:
                656 nm, 486 nm, 434 nm, \ldots\ in the visible range.)
                What does the mismatch suggest?
        \end{enumerate}

\end{enumerate}
\end{deepbox}

% ─── Reflection ───────────────────────────────────────────────────────────────
\subsection*{Reflection}

\begin{reflectionbox}
\begin{enumerate}[label=\textbf{R\arabic*.}]
  \item Thomson discovered the electron by studying something
        that nobody knew came from atoms (cathode rays).
        \textbf{What does this tell us about the role of
        ``curiosity-driven'' research ---
        studying phenomena without knowing their application ---
        versus directed research aimed at a specific goal?}

  \item The plum pudding model was a reasonable hypothesis
        given the evidence available in 1904.
        It was completely wrong about the distribution of
        positive charge.
        \textbf{What criterion should we use to judge a model:
        whether it is ultimately correct, or whether it was
        the best explanation of the available evidence?}

  \item Charge quantisation (Millikan's result) and mass
        conservation (Lavoisier's result) are both
        ``conservation-and-discreteness'' results.
        \textbf{What pattern do you notice about nature?}
        Make a general conjecture about other conserved,
        quantised quantities you might expect to find as
        we go deeper into physics.

  \item The electron is $1/1836$ the mass of a hydrogen atom.
        Atoms are mostly the mass of their nuclei.
        \textbf{What does this suggest about the relative
        importance of electrons vs nuclei for:
        (a) the mass of matter,
        (b) the chemical properties of matter,
        (c) the electrical properties of matter?}
\end{enumerate}
\end{reflectionbox}

\vspace{1.5em}
\begin{center}
\begin{tcolorbox}[
  enhanced, colback=DiscoveryColor!8, colframe=DiscoveryColor,
  width=0.88\textwidth, arc=3pt
]
\centering\sffamily\small
\textbf{What we invented in this module.}\\[4pt]
We invented the electron --- not as a philosophical speculation
but as a measured, quantified particle with a precise charge
and mass.
We also established the pattern that will recur throughout
Journey~2: every time we think we know what the atom looks like,
a single experiment shatters that picture.
The cathode ray tube showed Dalton's featureless sphere
was wrong: atoms have internal parts.
Now we must ask: how are those parts arranged?
Thomson's plum pudding model gives one answer.
Rutherford is about to destroy it.
\end{tcolorbox}
\end{center}

\vspace{2em}
\noindent\textcolor{JourneyColor!40}{\rule{\textwidth}{0.6pt}}
\begin{center}
\textit{\sffamily\small
Module 2.1 complete. Continue to Module 2.2:
\textbf{Rutherford's Gold Foil --- The Nuclear Surprise}.}
\end{center}

\end{document}